\documentclass[12pt, a4paper, oneside]{article}
\usepackage{amsmath, amsthm, amssymb, bm, graphicx, hyperref, mathrsfs, booktabs}
\usepackage[margin=2.5cm]{geometry}

\linespread{1.3}
\newtheorem{theorem}{Theorem}
\newtheorem{definition}[theorem]{Definition}
\newtheorem{remark}[theorem]{Remark}

\newcommand{\qvec}{\mathbf{q}}
\newcommand{\rvec}{\mathbf{r}}
\newcommand{\avec}{\mathbf{a}}
\newcommand{\bvec}{\mathbf{b}}
\newcommand{\qhat}{\hat{\mathbf{q}}}
\newcommand{\Nt}{N_t}
\newcommand{\Na}{N}
\newcommand{\Nq}{N_q}
\newcommand{\Nw}{N_\omega}
\newcommand{\field}{S}
\newcommand{\Wmat}{\mathbf{W}}

\begin{document}

\thispagestyle{empty}
\vspace*{\fill}
\begin{center}
    \Huge\textbf{Dynamic Structure Factor from Time-Correlation Functions}

    \vspace{1cm}
    \Large\textbf{Algorithm and Formulas}

    \vspace{0.8cm}
    \large \texttt{corr/sqw\_spin\_corr.py}
\end{center}
\vspace*{\fill}

\newpage
\tableofcontents
\newpage

% ======================================================================
\section{Overview}
% ======================================================================

Given a LAMMPS-like dump containing a three-component field $\field_j^\alpha(t)$
($\alpha = x,y,z$) on each atom $j = 1,\dots,\Na$ over $\Nt$ time frames, this
program computes the dynamic structure factor along a user-specified
$\qvec$-path through the time-correlation route:

\begin{equation}
    \field_j^\alpha(t) \;\longrightarrow\; s^\alpha(\qvec,t)
    \;\longrightarrow\; C^{ab}(\qvec,\tau)
    \;\longrightarrow\; S^{ab}(\qvec,\omega)
    \;\longrightarrow\; \text{channels}.
\end{equation}

The full $3\times3$ tensor is carried through every stage. This is what allows
a single run to emit several output channels at once, and --- with the FFT
evaluation of Sec.~\ref{sec:wk} --- it costs only a small multiple of a single
component rather than nine times as much.

\begin{remark}[Code layout]
The seven lines of \texttt{sqw\_at\_one\_q} in \texttt{sqw\_spin\_corr.py} are
the whole calculation, one per stage below. Parallelism appears nowhere in that
chain: the kernel is serial and handles one $\qvec$-point, and
\texttt{driver.py:map\_over\_q} is the only place that knows about MPI.
Supporting modules: \texttt{fields.py} (Sec.~\ref{sec:spatial}),
\texttt{correlate.py} (Sec.~\ref{sec:corr}), \texttt{spectrum.py}
(Sec.~\ref{sec:freq}), \texttt{channels.py} (Sec.~\ref{sec:channels}).
\end{remark}

% ======================================================================
\section{Input Data}
% ======================================================================

\subsection{Field Configuration}

The supercell contains $\Na$ atoms. The lattice is the $3\times3$ matrix of
row-vectors

\begin{equation}
    \mathbf{L} = \begin{pmatrix} \avec_1 \\ \avec_2 \\ \avec_3 \end{pmatrix},
\end{equation}

and at $t_n = n\Delta t$ we have the positions $\rvec_j(t_n)$ and the field
$\field_j^\alpha(t_n)$. The primitive cell is obtained by dividing $\mathbf{L}$
by the \texttt{-\/-supercell} repeats, and its reciprocal lattice, whose rows
satisfy $\bvec_i\cdot\avec_j^{\text{prim}} = 2\pi\delta_{ij}$, is

\begin{equation}
    \mathbf{B} = 2\pi\bigl(\mathbf{L}_{\text{prim}}^{-1}\bigr)^{\mathsf T}.
\end{equation}

All fractional $\qvec$ coordinates refer to this primitive reciprocal cell.

\subsection{Per-Atom Weights}

With \texttt{-\/-mass-weight}, each atom's field is scaled by $w_j =
\sqrt{m_{\text{type}(j)}}$ before the spatial sum (the phonon
spectral-energy-density convention). Without it $w_j = 1$, the magnon
convention.

\subsection{Field Threshold}

Atoms with $\max_t \lVert\field_j(t)\rVert \le \tau_{\text{spin}}$ are dropped
at load time to suppress noise from non-magnetic sites. Consequently $\Na$
denotes throughout this document the number of \emph{retained} atoms, and it is
that count which sets the $1/\sqrt{\Na}$ normalization below. The default
$\tau_{\text{spin}} = 0$ removes nothing.

% ======================================================================
\section{Spatial Fourier Transform}
\label{sec:spatial}
% ======================================================================

For each $\qvec$ (Cartesian, \AA$^{-1}$) and each frame,

\begin{equation}
\boxed{
    s^\alpha(\qvec,t) = \frac{\Lambda(\qvec)}{\sqrt{\Na}}
    \sum_{j=1}^{\Na} w_j\,\field_j^\alpha(t)\,e^{+i\qvec\cdot\rvec_j},
    \qquad \alpha = x,y,z.
}
\label{eq:spatial_ft}
\end{equation}

The $e^{+i\qvec\cdot\rvec}$ convention places a plane wave
$e^{i(\qvec\cdot\rvec-\omega t)}$ at the same signed $\qvec$ in the result, and
$1/\sqrt{\Na}$ makes $\lvert s\rvert^2$ intensive. Positions are frozen at the
first frame unless \texttt{-\/-use-instantaneous-pos} is given, in which case
$\rvec_j = \rvec_j(t)$.

Note that \emph{no projection is applied here}. The full Cartesian triple always
goes forward; which combination is reported is decided later by a weight matrix
(Sec.~\ref{sec:channels}).

\subsection{Finite-Size Lattice Sum}

\texttt{-\/-translation-repeats} $(N_1,N_2,N_3)$ repeats the loaded cell
explicitly, multiplying Eq.~\eqref{eq:spatial_ft} by

\begin{equation}
\boxed{
    \Lambda(\qvec) = \frac{1}{\sqrt{N_1N_2N_3}}
    \prod_{d=1}^{3}\sum_{n_d=0}^{N_d-1} e^{i n_d \qvec\cdot\avec_d},
}
\label{eq:lattice_sum}
\end{equation}

each axis sum being the geometric series $\bigl(1-e^{iN_d\theta_d}\bigr)/
\bigl(1-e^{i\theta_d}\bigr)$ with $\theta_d = \qvec\cdot\avec_d$. The prefactor
preserves the $1/\sqrt{\Na}$ normalization after repetition; $\Lambda = 1$ when
all repeats are $1$.

% ======================================================================
\section{Preparing the Time Series}
% ======================================================================

\subsection{Mean Subtraction}

Enabled by default, disabled by \texttt{-\/-no-subtract-mean}:

\begin{equation}
    \tilde s^\alpha(\qvec,t) = s^\alpha(\qvec,t) - \langle s^\alpha(\qvec)\rangle_t,
    \qquad
    \langle s^\alpha\rangle_t = \frac{1}{\Nt}\sum_{t} s^\alpha(\qvec,t),
\end{equation}

which removes the elastic ($\omega = 0$) line and any static background. It must
precede the window: tapering a series that still carries a DC offset would
imprint the shape of the window itself onto the spectrum.

\subsection{Data Window}
\label{sec:window}

\begin{equation}
\boxed{
    \tilde s^\alpha_{\text{win}}(\qvec,t) = w(t)\,\tilde s^\alpha(\qvec,t),
    \qquad t = 0,\dots,\Nt-1,
}
\label{eq:window}
\end{equation}

with $w$ either the Hann taper $\tfrac12\bigl[1-\cos(2\pi n/(\Nt-1))\bigr]$
(\texttt{-\/-window hann}, default) or unity (\texttt{-\/-window none}).

\begin{remark}[Why a data window and not a lag window]
The window acts on $s(\qvec,t)$, \emph{before} the correlation is formed. Its
job is to suppress \textbf{spectral leakage}. The record is finite and the
transform treats it as periodic, so unless a mode completes a whole number of
cycles within the record the splice between last and first frame is a
discontinuity, and its power smears across every frequency. Pulling both ends to
zero removes that discontinuity. Measured on a two-mode test signal, a weak
branch $10^{-6}$ below a strong one is overestimated by a factor $\sim 800$
without the window and by a factor $1.1$ with it.

Tapering the correlation $C(\tau)$ instead --- a lag window, the
Blackman--Tukey construction --- addresses a different problem, the statistical
scatter of the estimate. It is deliberately \textbf{not} offered here, because
it is only usable when the physical linewidth greatly exceeds the frequency
resolution $1/(\Nt\Delta t)$. Below that threshold it inflates the measured
linewidth by an order of magnitude, destroying the very quantity a
damped-oscillator fit is after. When variance genuinely needs reducing, the
correct remedy is to average independent trajectories or symmetry-equivalent
$\qvec$-points; lengthening a single trajectory does not help, since the extra
frames buy a finer frequency grid rather than a better value per bin.
\end{remark}

% ======================================================================
\section{Time-Correlation Tensor}
\label{sec:corr}
% ======================================================================

\begin{equation}
\boxed{
    C^{ab}(\qvec,\tau) = \frac{1}{M_\tau}\sum_{t=0}^{\Nt-\tau-1}
    \tilde s^a_{\text{win}}(\qvec,t+\tau)\,
    \tilde s^{b*}_{\text{win}}(\qvec,t),
    \qquad a,b\in\{x,y,z\},
}
\label{eq:corr}
\end{equation}

for $\tau = 0,\dots,\Nt-1$, with

\begin{itemize}
    \item \textbf{biased} (\texttt{-\/-corr-norm biased}, default):
    $M_\tau = \Nt$. Constant denominator; keeps the estimate a genuine power
    spectrum, i.e.\ positive semidefinite.
    \item \textbf{unbiased} (\texttt{-\/-corr-norm unbiased}):
    $M_\tau = \Nt-\tau$. Corrects for the shrinking number of contributing
    pairs, but destroys positive semidefiniteness, so the transform can go
    negative. See Remark~\ref{rem:corrplus} for when this is nevertheless the
    right choice.
\end{itemize}

\subsection{Evaluation by FFT}
\label{sec:wk}

Equation~\eqref{eq:corr} is not evaluated as written. By the Wiener--Khinchin
theorem,

\begin{equation}
\boxed{
    C^{ab} \;\propto\; \mathcal{F}^{-1}\Bigl[\,
        \hat s^a \cdot \overline{\hat s^b}\,\Bigr],
    \qquad \hat s^a = \mathcal{F}\bigl[\tilde s^a_{\text{win}}\bigr],
}
\label{eq:wk}
\end{equation}

where the transform is taken over a sequence zero-padded to length
$L \ge 2\Nt$. The padding makes the circular correlation of the padded signal
equal the linear correlation of the original, so Eq.~\eqref{eq:wk} is an exact
identity rather than an approximation.

Three forward transforms (one per Cartesian component) and nine inverse
transforms give the whole tensor, at $O(\Nt\log\Nt)$ instead of the $O(\Nt^2)$
of the explicit double loop. Measured speed-ups range from $\sim50\times$ at
$\Nt = 2\times10^3$ to $\sim450\times$ at $\Nt = 2\times10^4$; the gap exceeds
the raw operation count because the direct form is a Python-level loop over
$\Nt$ separate array calls.

All of this runs in complex128 regardless of the storage dtype of the
trajectory. Since $s(\qvec,t)$ is only $(\Nt,3)$ the promotion is free in
memory, while the long accumulations would otherwise eat well into float32's
seven digits.

\subsection{Extension to Negative Lags}

\begin{equation}
\boxed{
    C^{ab}(\qvec,-\tau) = \overline{C^{ba}(\qvec,\tau)},
    \qquad \tau = 1,\dots,\Nt-1,
}
\label{eq:herm}
\end{equation}

giving a two-sided sequence of length $2\Nt-1$ indexed
$\tau = -(\Nt-1),\dots,\Nt-1$. Note the \textbf{index swap}: writing
$\overline{C^{ab}(\tau)}$ instead would silently give the wrong answer for every
off-diagonal element. This symmetry is what makes the following transform real.

% ======================================================================
\section{Frequency Domain}
\label{sec:freq}
% ======================================================================

\begin{equation}
\boxed{
    S^{ab}(\qvec,\omega_k) = \sum_{\tau=-(\Nt-1)}^{\Nt-1}
    e^{-i\omega_k\tau\Delta t}\,C^{ab}(\qvec,\tau),
    \qquad \omega_k \ge 0,
}
\label{eq:sqw}
\end{equation}

evaluated by FFT after \texttt{np.fft.ifftshift} maps the two-sided sequence
into FFT order. There is \textbf{no $\Delta t$ prefactor}: the code returns the
bare DFT sum, so $S$ is a relative intensity and absolute values are not
comparable across runs that use different $\Delta t$.

At each frequency the result is a $3\times3$ matrix, Hermitian in $(a,b)$ by
Eq.~\eqref{eq:herm} and positive semidefinite as a spectral-density matrix ---
the property Sec.~\ref{sec:channels} relies on.

The frequency grid follows the transform length $L = 2\Nt-1$:

\begin{equation}
    f_k = \frac{k}{(2\Nt-1)\,\Delta t\,[\mathrm{s}]\times10^{12}}\ \mathrm{THz},
    \qquad k = 0,\dots,\Nt-1.
\end{equation}

Because $L$ is odd there are exactly $\Nt$ non-negative bins, the largest being
$(\Nt-1)/[(2\Nt-1)\Delta t]$, which approaches but never reaches the Nyquist
frequency $1/(2\Delta t)$.

\begin{remark}[Why only $\omega \ge 0$]
Discarding negative frequencies is not a consequence of Eq.~\eqref{eq:herm}.
That makes $S$ real but \emph{not} even in $\omega$: a single propagating mode
$s(\qvec,t) = Ae^{-i\omega_0 t}$ gives $C(\tau) = \lvert A\rvert^2
e^{-i\omega_0\tau}$, which peaks at $\omega = -\omega_0$ only. Evenness holds
when the dynamics are classical \emph{and} the system is centrosymmetric:
reality of the field gives $S(\qvec,\omega) = S(-\qvec,-\omega)$ and inversion
symmetry gives $S(\qvec,\omega) = S(-\qvec,\omega)$. For chiral systems
(Dzyaloshinskii--Moriya, handed magnons) that fails and weight residing at
$\omega < 0$ is genuinely lost.
\end{remark}

\begin{remark}[Relation to the periodogram]
With the data window of Sec.~\ref{sec:window}, $C$ is the autocorrelation of the
windowed series, so by Eq.~\eqref{eq:wk} its transform is exactly
$\lvert\mathcal{F}[\tilde s\cdot w]\rvert^2$ --- a zero-padded windowed
periodogram. The two routes agree to $6\times10^{-16}$ for both window choices.
The only difference is the frequency grid: $2\Nt-1$ bins here against $\Nt$ for
a direct periodogram, and the extra bins are zero-padding interpolation, not
improved resolution, which remains set by $\Nt\Delta t$.

What the correlation route retains that a periodogram does not is direct access
to $C^{ab}(\qvec,\tau)$ itself (Remark~\ref{rem:corrplus}) and the
\texttt{unbiased} normalization.
\end{remark}

% ======================================================================
\section{Output Channels}
\label{sec:channels}
% ======================================================================

Every quantity this program can report is a real linear form on the tensor:

\begin{equation}
\boxed{
    \text{channel}(\qvec,\omega)
    = \sum_{a,b} W_{ab}\,\operatorname{Re}S^{ab}(\qvec,\omega),
}
\label{eq:channel}
\end{equation}

so a channel is completely specified by one real $3\times3$ weight matrix
$\Wmat$. The \texttt{-\/-component} tokens are names for such matrices.

\subsection{Token Grammar}

\begin{center}
\texttt{token := term ('+' term)*}
\end{center}

Tokens separated by spaces produce separate outputs; terms joined by \texttt{+}
are summed into one. Thus \texttt{-\/-component 1+2 3} emits two curves,
$S^{xx}+S^{xy}$ and $S^{xz}$.

\begin{table}[h]
    \centering
    \begin{tabular}{lll}
        \toprule
        Term & $\Wmat$ & Meaning \\
        \midrule
        \texttt{1}\,\dots\,\texttt{9} & $e_ae_b^{\mathsf T}$ & row-major: $1=xx$, $2=xy$, $3=xz$, $4=yx$, \dots \\
        \texttt{xx}\,\dots\,\texttt{zz} & $e_ae_b^{\mathsf T}$ & same, named by axis pair \\
        \texttt{x}, \texttt{y}, \texttt{z} & $e_ae_a^{\mathsf T}$ & shorthand for \texttt{xx}, \texttt{yy}, \texttt{zz} \\
        \texttt{L} & $\qhat\qhat^{\mathsf T}$ & longitudinal to $\qvec$ \\
        \texttt{T} & $\mathbf{I}-\qhat\qhat^{\mathsf T}$ & transverse to $\qvec$ \\
        \bottomrule
    \end{tabular}
    \caption{Terms and their weight matrices. \texttt{L} and \texttt{T} depend
    on the direction of $\qvec$ and are rebuilt per $\qvec$-point; the rest are
    constants. The default \texttt{1+5+9} is the trace.}
\end{table}

At $\qvec\to0$ the longitudinal direction is undefined: \texttt{L} contributes
nothing and \texttt{T} reduces to the full trace, the natural $\Gamma$-point
limit.

\subsection{Equivalence with Projections}

A projection of the field is a special case of Eq.~\eqref{eq:channel}. For the
longitudinal case, $s^L = \sum_a\hat q^a s^a$ with real $\qhat$ gives

\begin{equation}
    C^L(\tau) = \sum_{a,b}\hat q^a\hat q^b\,C^{ab}(\tau),
\end{equation}

which is exactly the \texttt{L} weight. For the transverse case,
$P = \mathbf{I}-\qhat\qhat^{\mathsf T}$ is real, symmetric and idempotent, so

\begin{equation}
    \sum_a C^{T,aa}(\tau)
    = \sum_{c,d}\bigl(P^{\mathsf T}P\bigr)_{cd}C^{cd}(\tau)
    = \sum_{c,d}P_{cd}\,C^{cd}(\tau),
\end{equation}

the \texttt{T} weight. Both identities were confirmed numerically against an
independent projection implementation, agreeing to $\sim10^{-8}$, the precision
of complex64 arithmetic.

\subsection{Clipping}

$S^{ab}$ is Hermitian positive semidefinite at each frequency, so

\begin{equation}
    \sum_{a,b}W_{ab}S^{ab} \ge 0
    \quad\Longleftrightarrow\quad
    \Wmat \text{ is symmetric positive semidefinite}.
\end{equation}

That is precisely the condition under which the code clips negative values ---
which are then numerical noise --- to zero. Channels whose $\Wmat$ is not
symmetric PSD are signed quantities and keep their sign: any single off-diagonal
element, and mixed groups such as \texttt{1+2}. This is the exact generalisation
of the older diagonal/off-diagonal rule, which cannot classify a group.

\begin{remark}[The imaginary part]
For $u \ne v$ the quantity $S^{uv}$ is complex, and Eq.~\eqref{eq:channel}
retains only its real part. The antisymmetric imaginary part carries the chiral
information (spin chirality, Dzyaloshinskii--Moriya terms) and is currently
discarded.
\end{remark}

% ======================================================================
\section{Output}
% ======================================================================

\subsection{NPZ Contents}

\begin{itemize}
    \item \texttt{timesteps} --- LAMMPS timestep of each retained frame,
    \item \texttt{q\_frac}, \texttt{q\_vectors} --- $\qvec$ in primitive fractional and Cartesian (\AA$^{-1}$) coordinates,
    \item \texttt{q\_path\_distance}, \texttt{q\_node\_indices} --- plot abscissa and high-symmetry node positions,
    \item \texttt{freq\_thz}, \texttt{energy\_meV} --- frequency grid, in THz and in meV ($E = hf$, $h \approx 4.13567$~meV$\cdot$ps),
    \item \texttt{sqw} --- the channel, shape $(\Nq,\Nw)$,
    \item \texttt{components} --- the token that produced it,
    \item \texttt{field\_columns}, \texttt{lattice}, \texttt{reciprocal\_lattice}, \texttt{translation\_repeats}, \texttt{dt\_fs}, \texttt{corr\_norm} --- provenance,
    \item \texttt{corr\_plus} --- present only with \texttt{-\/-save-corr-plus}; the one-sided $C^{ab}(\qvec,\tau)$, shape $(\Nq,\Nt,3,3)$.
\end{itemize}

A single requested channel writes to \texttt{-\/-output} unchanged; several
channels each get their own file with a \texttt{\_<token>} suffix, with
\texttt{+} rendered as \texttt{p}.

\subsection{The Correlation Itself}

\begin{remark}[Reading a lifetime from \texttt{-\/-save-corr-plus}]
\label{rem:corrplus}
$C(\qvec,\tau)$ is the time-domain view of the same physics: its oscillation
period is the mode frequency and the decay of its envelope is the linewidth,
i.e.\ the magnon lifetime. Fitting that envelope is often more robust than
fitting a Lorentzian in frequency, because the frequency estimate scatters by
$\sim100\%$ per bin while $C(\tau)$ at small $\tau$ is determined by many sample
pairs.

Two settings will corrupt the envelope, and both are silent:
\begin{itemize}
    \item \textbf{biased normalization} divides by a constant $\Nt$ while only
    $\Nt-\tau$ terms contribute, imprinting a spurious triangular decay
    $(\Nt-\tau)/\Nt$. Use \texttt{-\/-corr-norm unbiased}, which cancels it
    exactly.
    \item \textbf{the data window} multiplies $C$ by the window's own
    autocorrelation. Use \texttt{-\/-window none}.
\end{itemize}
On an undamped coherent test signal, whose true envelope is flat, the default
\texttt{hann}/\texttt{biased} combination decays to $0.7\%$ by $3$~ps --- a
lifetime that is entirely an artefact of the analysis. The combination
\texttt{-\/-window none -\/-corr-norm unbiased} reproduces the flat envelope.

The frequency-domain and time-domain uses therefore want \emph{opposite}
settings: \texttt{hann}/\texttt{biased} for $S(\qvec,\omega)$,
\texttt{none}/\texttt{unbiased} for $C(\qvec,\tau)$.
\end{remark}

% ======================================================================
\section{Summary of the Algorithm}
% ======================================================================

For each $\qvec$-point on the path:

\begin{enumerate}
    \item \textbf{Spatial FT} --- Eq.~\eqref{eq:spatial_ft}, including the
    lattice-sum factor $\Lambda(\qvec)$ of Eq.~\eqref{eq:lattice_sum}.
    \item \textbf{Subtract the mean} --- removes the elastic line.
    \item \textbf{Data window} --- Eq.~\eqref{eq:window}, suppresses leakage.
    \item \textbf{Correlate} --- the full tensor Eq.~\eqref{eq:corr}, evaluated
    by the zero-padded FFT identity Eq.~\eqref{eq:wk}.
    \item \textbf{Extend to negative lags} --- Eq.~\eqref{eq:herm}, minding the
    index swap.
    \item \textbf{Transform} --- Eq.~\eqref{eq:sqw}, keeping $\omega \ge 0$.
    \item \textbf{Contract with each channel weight} --- Eq.~\eqref{eq:channel},
    clipping only where $\Wmat$ is symmetric PSD.
\end{enumerate}

Steps 1--7 form the serial per-$\qvec$ kernel; \texttt{map\_over\_q}
distributes it across MPI ranks and gathers the result. Because each
$\qvec$-point is independent, results are bitwise identical for any rank count.

% ======================================================================
\section{Notation}
% ======================================================================

\begin{table}[h]
    \centering
    \begin{tabular}{ll}
        \toprule
        Symbol & Meaning \\
        \midrule
        $\Na$ & Atoms retained after the field threshold \\
        $\Nt$ & Time frames \\
        $\Nq$, $\Nw$ & $\qvec$-points, frequency bins \\
        $\Delta t$ & Interval between saved frames (fs) \\
        $\qvec$, $\qhat$ & Wave-vector (\AA$^{-1}$) and its unit vector \\
        $\field_j^\alpha(t)$ & Field component $\alpha$ on atom $j$ \\
        $w_j$, $w(t)$ & Per-atom mass weight; data window \\
        $s^\alpha(\qvec,t)$ & Spatial Fourier amplitude, Eq.~\eqref{eq:spatial_ft} \\
        $\Lambda(\qvec)$ & Lattice-sum factor, Eq.~\eqref{eq:lattice_sum} \\
        $C^{ab}(\qvec,\tau)$ & Correlation tensor, Eq.~\eqref{eq:corr} \\
        $S^{ab}(\qvec,\omega)$ & Structure-factor tensor, Eq.~\eqref{eq:sqw} \\
        $\Wmat$ & Channel weight matrix, Eq.~\eqref{eq:channel} \\
        $M_\tau$ & Correlation denominator: $\Nt$ or $\Nt-\tau$ \\
        \bottomrule
    \end{tabular}
    \caption{Notation.}
\end{table}

\end{document}
